\begin{helper}[GnpAllVerticesDegreeFailureUnionBound]
Let \(m\in\mathbb N\), \(0\le p\le1\), \(\epsilon,t\in\mathbb R\), and
\(0\le t\).  Under the \(\noderef{GnpGraphWeight}\) product graph weight,
the mass of the event that some vertex has degree outside
\(\noderef{WithinRelativeError}(\epsilon,pm,\cdot)\) is at most
\[
m\left[
\exp\!\bigl(-t(1+\epsilon)pm\bigr)
\bigl(1-p+p\exp(t)\bigr)^{m-1}
+
\exp\!\bigl(t(1-\epsilon)pm\bigr)
\bigl(1-p+p\exp(-t)\bigr)^{m-1}
\right].
\]
\end{helper}
\begin{proof}
For a fixed vertex \(v\), failure of
\(\noderef{WithinRelativeError}(\epsilon,pm,\noderef{GraphDegreeCount}(G,v))\)
means
\[
|\noderef{GraphDegreeCount}(G,v)-pm|>\epsilon pm.
\]
Therefore either
\(\noderef{GraphDegreeCount}(G,v)\ge (1+\epsilon)pm\) or
\(\noderef{GraphDegreeCount}(G,v)\le (1-\epsilon)pm\).  The graph weights are
nonnegative because \(0\le p\le1\), so the finite union bound over
\(v\in\operatorname{Fin}(m)\) bounds the all-vertices failure mass by the sum
over vertices of these two fixed-vertex tail masses.

For each \(v\), the number of incident edge coordinates is \(m-1\): it is in
bijection with the vertices other than \(v\).  Applying
\(\noderef{GnpVertexDegreeUpperTailUnionBound}\) to the upper and lower tails
at this vertex therefore gives the displayed two-tail bound.  Summing the
same bound over all \(m\) vertices gives the stated factor of \(m\).
\end{proof}
